\newpage


%%%%%%%%%%%%%%%%%%%%%%%%%%%%%%%%%%%%%%%%%%%%%%%%%%
%%%%%%%%%%%%%%%%%%%%%%%%%%%%%%%%%%%%%%%%%%%%%%%%%%
\section{The Problem}
%%%%%%%%%%%%%%%%%%%%%%%%%%%%%%%%%%%%%%%%%%%%%%%%%%
%%%%%%%%%%%%%%%%%%%%%%%%%%%%%%%%%%%%%%%%%%%%%%%%%%

Let $\Q$ be a $\Q$-like set (i.e. dense,ordered abelian group with archemedian property that contains more than just the zero element --Definition~\ref{def:qlike}); We consider $\Q$ with its natural metric
\[ d: \Q \times \Q \to \Q \mid d(x,y) = |x-y|, \]
whose topology is the order topology (Lemma~\ref{lem:natmetric}). Our problem
is that this metric need not be complete: for the rationals it has holes
(Theorem~\ref{thm:holes}). We therefore wish to find a larger metric space
$\Qb$ such that
\[
\begin{aligned}
(1)\quad & (\Q,d) \subseteq (\Qb,\db),\\[2mm]
(2)\quad & (\Qb,\db)\ \text{is complete},\\[2mm]
(3)\quad & \text{if } Z \text{ is any complete metric space with }
(\Q,d)\subseteq (Z,d_Z),\text{ then } (\Qb,\db)\subseteq (Z,d_Z).
\end{aligned}
\]
These properties make $\Qb$ the set containing $\Q$ (1), all of the missing
holes (2), and nothing else (3). Here ``$\subseteq$'' means ``embeds
isometrically into''; this is made precise in Section~\ref{sec:embed} and
Section~\ref{sec:smallest}. I now construct such a set and prove it has the
desired properties.


%%%%%%%%%%%%%%%%
%%%%%%%%%%%%%%%%
%%%%%%%%%%%%%%%%
% formal construction
%%%%%%%%%%%%%%%%
%%%%%%%%%%%%%%%%
%%%%%%%%%%%%%%%%

\subsection{Formal Construction of $\Qb$} \label{sec:construct}

The plan: construct $\Qb$ as equivalence classes of Cauchy sequences
(Step 1); construct $\db:\Qb\times\Qb\to\Qb$ and show it is well defined
(Step 2); show that $\Qb$ is itself a valid set of values, i.e.\ a $\Q$-like
set (Step 3); and show that $\db$ is a metric, in fact the natural one
(Step 4).

\step{1}{Construct the ambient space.}

Write $\mathscr{L}(\N \to \Q)$ for the set of all functions $\N\to\Q$, i.e.\
all sequences in $\Q$. It exists by the power set axiom: a function is a
subset of $\N\times\Q$, so $\mathscr{L}(\N\to\Q)$ is carved out of
$\mathcal P(\N\times\Q)$ by specification. By specification again we keep only
the Cauchy sequences:
\[
S := \{(x_n)\in\mathscr{L}(\N \to \Q) \mid (x_n) \text{ is a Cauchy sequence in } (\Q,d)\}.
\]
Define the relation on $S$:
\[
(x_n) \sim (y_n)
\iff d\text{-}\lim_{\substack{n \to \infty \\ n \in \N }} d(x_n, y_n) = 0 .
\]

\begin{lemma} $\sim$ is an equivalence relation on $S$.
\end{lemma}

This defines equivalence classes:
\[
[(x_n)] = \{(z_n) \in S \mid (z_n) \sim (x_n)\},
\qquad
[(x_n)] = [(y_n)] \iff (x_n) \sim (y_n).
\]
By the axiom of replacement (or by specification inside $\mathcal P(S)$), we
replace each $(x_n) \in S$ with its equivalence class $[(x_n)]$. Thus the
ambient space is:
\[
\Qb :=
\{[(x_n)] \mid (x_n) \in \mathscr{L}(\N \to \Q),\ (x_n)\ \text{Cauchy}\}.
\]

%The idea: a Cauchy sequence ``points at'' a location, possibly a hole, and two
%sequences are identified when they point at the same location.
% rather not have this sentence, keep commented out

\step{2}{Define the distance.}

\[ \db: \Qb \times \Qb \to \Qb, \qquad
 \db( [(x_n)] ,[(y_n)] )  :=  [(d(x_n,y_n))] . \]

I must show that (a) $[(d(x_n,y_n))] $ is indeed in $\Qb$ and (b) that this is a well defined function ...

\step{3}{$\Qb$ is a $\Q$-like set.}

Before we can even say that $\db$ is a metric we need ``$+$'' and ``$<$'' on
$\Qb$, satisfying (V1) and (V2). We also check (V3), density and
nontriviality, so that $\Qb$ is as good a set of values as $\Q$ was.

\substep{a}{\textbf{(V1) abelian group.} Define}
\[
[(x_n)]+[(y_n)] := [(x_n+y_n)],\qquad -[(x_n)] := [(-x_n)],\qquad 0:=[(0,0,0,\dots)].
\]

The abelian group laws hold because they hold term by term in $\Q$. For
$a\in\Q$ we write $\iota(a):=[(a,a,a,\dots)]$ for the class of the constant
sequence; I come back to this in section ~\ref{sec:embed}. Note
$\iota(a)+\iota(b)=\iota(a+b)$ and $-\iota(a)=\iota(-a)$.

\substep{b}{\textbf{(V2) positivity and order.}}

\begin{definition} (positivity, order) $\mid$
Call $[(x_n)]$ \emph{positive}, written $[(x_n)]>0$, if
\[
\exists \delta\in\Q,\ \delta>0,\ \exists N \mid \forall n>N,\ x_n\ge \delta .
\]
Define $X<Y \iff Y-X>0$, and $X\le Y\iff X<Y$ or $X=Y$.
\end{definition}

% Positivity is well defined, and $\iota$ preserves order:
% \begin{align*}
% &(y_n)\sim(x_n),\ x_n\ge\delta\ev,\ \eta:=\delta/2
% \imp |x_n-y_n|<\eta\ev \imp y_n>x_n-\eta\ge\delta-\eta\ge\eta>0\ev\\
% &\qquad\imp [(y_n)]>0.\\
% &a<b \imp b-a\ge\delta:=b-a>0 \imp \iota(b)-\iota(a)=\iota(b-a)>0;\qquad
% \iota(a)<\iota(b) \imp b-a\ge\delta>0 \imp a<b.
% \end{align*}

\begin{lemma} (trichotomy) $\mid$ \label{lem:tri}
For every $X=[(x_n)]\in\Qb$ exactly one of $X=0$, $X>0$, $-X>0$ holds.
Consequently $\Qb$ satisfies (V1) and (V2).
\end{lemma}

The following lemma is the bridge between orderings of sequences and
orderings of classes; we use it constantly.

\begin{lemma} (eventual comparison) $\mid$ \label{lem:evcomp}
For $(u_n),(v_n)\in S$:
\[
\exists N \mid \forall n>N, u_n\le v_n \imp [(u_n)]\le[(v_n)] .
\]
\end{lemma}

\begin{lemma} ($\Q$ is coinitial and cofinal in $\Qb$) $\mid$ \label{lem:ratbounds}
\begin{enumerate}[nosep,leftmargin=3em]
    \item[(a)] $\epsilon\in\Qb,\ \epsilon>0 \imp \exists q\in\Q:\ 0<\iota(q)<\epsilon$.
    \item[(b)] $X\in\Qb \imp \exists q\in\Q:\ X<\iota(q)$.
\end{enumerate}
\end{lemma}

\substep{c}{\textbf{(V3) Archimedean.}}
\begin{align*}
&X,Y>0 \imp x_n\ge\delta>0\ev \imp X\ge\iota(\delta);\qquad
Y<\iota(M)\ \text{ for some } M\in\Q\ \ (\ref{lem:ratbounds}\text{b});\\
&\text{(V3) in } \Q \imp \exists n:\ n\delta>M \imp nX\ge n\,\iota(\delta)=\iota(n\delta)>\iota(M)>Y.
\end{align*}

\substep{d}{\textbf{density, nontriviality.}}
\begin{align*}
&X<Y \imp Y-X>0 \by{\ref{lem:ratbounds}a} \exists q\in\Q:\ 0<\iota(q)<Y-X
\imp X<X+\iota(q)<Y.\\
&c>0 \text{ in } \Q \imp \iota(c)>0 \imp \Qb\ne\{0\}.
\end{align*}

\step{4}{$(\Qb,\db)$ is a metric space with values in $\Qb$, and $\db$ is the natural metric of $\Qb$.}

\begin{lemma} ($\db$ is the natural metric of $\Qb$) $\mid$ \label{lem:natbar}
For all $X,Y\in\Qb$: $\ \db(X,Y)=|Y-X|$, with $|\cdot|$ taken in the
$\Q$-like set $\Qb$. So $\Qb$ carries exactly the same kind of metric as $\Q$
did (Definition~\ref{def:natmetric}), and its topology is the order topology
(Lemma~\ref{lem:natmetric}).
\end{lemma}


